\headright{Expérience}

\jobtitle{Ingénieur DevOps}{Alten (France)}{mai 2025 -- présente}{}
\small

\noindent \textbf{Amadeus}

\begin{itemize}[label=\tiny$\bullet$, leftmargin=1em]
	\item Migration des scripts de pipeline Groovy hérités vers un framework interne gérant l'intégralité du cycle de vie des applications.
	\item Réécriture du framework de flux de données de \textbf{Bash} vers \textbf{Python} pour les livraisons de données planifiées.
	\item Mise en place de frameworks industriels pour automatiser les dépendances locales (Localstack, Nix devenv).
	\item Développement de scripts permettant de basculer dynamiquement les types de machines sur les clusters EMR de spot vers on-demand en fonction de la latence de la file d'attente.
	\item Refonte de la documentation interne pour améliorer la clarté et la maintenabilité.
	\item Conception d'un HLD pour la gestion des versions, standardisant les workflows de déploiement conformément aux exigences de gestion des changements.
\end{itemize}

\vspace{3.0ex}
\jobtitle{Ingénieur DevOps}{DXC Technology (Bulgarie)}{02/2022 -- 04/2025} {}
\small

\noindent \textbf{Continental AG (Pneus)}
\begin{itemize}[label=\tiny$\bullet$, leftmargin=1em, nosep]
	\item Direction d'une équipe de 5 personnes et compte rendu direct de l'état d'avancement du projet au client.
	\item Optimisation de la logique de traitement des données par la migration des tâches batch Spark héritées vers des pipelines de streaming Apache Flink, réduisant ainsi considérablement la latence.
	\item Accélération du lancement du projet grâce à la résolution de problèmes critiques de connectivité VPN et à l'automatisation des procédures de développement local.
\end{itemize}

\vspace{1em}
\noindent \textbf{Nestlé}
\begin{itemize}[label=\tiny$\bullet$, leftmargin=1em, nosep]
	\item Mise en place de contrôles automatisés de validation structurelle pour les Pull Requests afin d'empêcher les déploiements non conformes.
	\item Extension des dialectes Snowflake dans SQLFluff et intégration automatisée lors de l'étape de pré-lancement
	\item Intégration des analyseurs syntaxiques ANTLR dans le pipeline \textbf{CI/CD} en tant que porte de validation obligatoire avant le lancement.
\end{itemize}

\vspace{1em}
\noindent \textbf{Ahold Delhaize}
\begin{itemize}[label=\tiny$\bullet$, leftmargin=1em]
	\item Refonte du cadre d'automatisation des tests en une architecture modulaire, améliorant la réutilisabilité du code et réduisant le temps d'intégration pour les ingénieurs QA.
	\item Collaboration avec l'équipe QA lors de réunions hebdomadaires pour identifier les goulots d'étranglement opérationnels, ce qui a permis d'améliorer l'autonomie de l'équipe.
	\item Conception d'un HLD pour le partage de données entre environnements, qui a permis la migration du nouveau cadre d'ingestion cloud (CF2).
	\item Automatisation d'un flux pour héberger les rapports Allure sous forme de site web statique sur Azure Blob Storage avec un middleware d'authentification personnalisé.
	\item Direction d'une évaluation technique de \textbf{SonarQube} pour l'analyse automatisée du code, afin d'identifier les exigences d'intégration pour la plateforme
\end{itemize}

\vspace{1em}
\noindent \textbf{Continental AG (ADAS, nom de code CAEdge)}
\begin{itemize}[label=\tiny$\bullet$, leftmargin=1em]
	\item Automatisation d'une topologie \textbf{VPC} standardisée pour une configuration AWS multi-comptes, en l'intégrant à l'orchestrateur Account Vending Machine.
	\item Conception et automatisation d'une solution de gouvernance des coûts cloud au sein d'une organisation AWS multi-comptes, avec mise en place d'alertes budgétaires pour prévenir les dépassements.
	\item Réduction de la surface d'attaque des applications cloud grâce à la correction des vulnérabilités (CVE) identifiées par les rapports AWS Inspector.
\end{itemize}

\vspace{5.0ex}
\jobtitle{Ingénieur logiciel stagiaire}{Sys Consulting Ltd.}{07/2021 -- 12/2021}{}
\small
Maintenance de systèmes bancaires hérités dans CAVO pour des environnements Windows XP.
